\documentclass{article}
\usepackage{amsmath, amssymb, amsthm}
\usepackage{hyperref}
\usepackage{geometry}
\geometry{margin=1in}

\title{Erd\H{o}s Problem \#1200}
\date{Last updated: 2026-04-04}
\author{Source: erdosproblems.com}

\begin{document}
\maketitle

\noindent\textbf{Status:} OPEN \quad \textbf{No prize}

\noindent\textbf{Tags:} number theory, primes

\medskip

\noindent\textbf{Problem Statement:}

There exists a constant $C$ such that for all large $x$ there is a collection of primes $p_1&#60;\ldots&#60;p_k&#60;x$ with $\sum\frac{1}{p_i}&#60;C$ together with a system of congruences $a_i\pmod{p_i}$ such that every integer $n&#60;x$ satisfies at least one of these congruences.




A conjecture of Erd\H{o}s and Ruzsa which Erd\H{o}s described in \cite{Er80} as 'surprising'. He says that if this is correct then 'very likely' the $\epsilon_n$ in [688] is $\gg 1$ - certainly proving $\epsilon_n\geq c$ would prove this conjecture (taking $P$ to be all primes in $[x^c,x]$). In \cite{ErRu80} this is asked as a question: if $p_1&lt;\cdots&lt;p_k&lt;x$ are primes with $\sum \frac{1}{p_i}\leq C$ and $a_i\pmod{p_i}$ are any residue classes then must there always be $\gg_K x$ many integers $n&lt;x$ avoiding all of them? (Of course is the answer is yes then this disproves the main conjecture.) Erd\H{o}s and Ruzsa \cite{ErRu80} proved that for any $C&gt;0$ there exists a set of primes $P$ such that $\sum_{p\in P}\frac{1}{p}\leq C$ and the number of integers $n\leq x$ divisible by at least one $p\in P$ is $\gg_C x$. See also [783] and [784] .




References

[Er80] Erd\H{o}s, Paul, A survey of problems in combinatorial number theory . Ann. Discrete Math. (1980), 89-115.

[ErRu80] Erd\H{o}s, P. and Ruzsa, I. Z., On the small sieve. I. Sifting by primes . J. Number Theory (1980), 385--394.

\noindent\textbf{Related OEIS sequences:} possible


\bigskip
\noindent\small{Source: \url{https://www.erdosproblems.com/1200}}
\end{document}
